
\documentclass{article}
\usepackage{colortbl}
\usepackage{makecell}
\usepackage{multirow}
\usepackage{supertabular}

\begin{document}

\newcounter{utterance}

\centering \large Interaction Transcript for game `ifeval', experiment `ifeval', episode 429 with Qwen3.5{-}4B.
\vspace{24pt}

{ \footnotesize  \setcounter{utterance}{1}
\setlength{\tabcolsep}{0pt}
\begin{supertabular}{c@{$\;$}|p{.15\linewidth}@{}p{.15\linewidth}p{.15\linewidth}p{.15\linewidth}p{.15\linewidth}p{.15\linewidth}}
    \# & \multicolumn{2}{c}{Player} && \multicolumn{2}{c}{Game Master} \\
    \hline

    \theutterance \stepcounter{utterance}  
    & & & \multicolumn{4}{p{0.6\linewidth}}{
        \cellcolor[rgb]{0.9,0.9,0.9}{
            \makecell[{{p{\linewidth}}}]{
                \texttt{\tiny{[P1$\langle$GM]}}
                \texttt{Is "jiang" a Chinese name? What are other names similar to "jiang"? Separate your two answers with ******} \\
            }
        }
    }
     \\ \\

    \theutterance \stepcounter{utterance}  
    & \multicolumn{4}{p{0.6\linewidth}}{
        \cellcolor[rgb]{0.561,0.737,0.561}{
            \makecell[{{p{\linewidth}}}]{
                \texttt{\tiny{[P1$\rangle$GM]}}
                \texttt{Yes, "Jiang" (江) is a very common Chinese surname, ranking as the 10th most frequent surname in China. It literally means "river" and is often associated with the Yangtze River, which flows through the region where many people with this surname live.} \\
\texttt{******} \\
\texttt{Other surnames similar to "Jiang" in terms of meaning (relating to water, rivers, or flowing) or sound include:} \\
\texttt{*   **Song** (宋): Often associated with the Song River.} \\
\texttt{*   **Shen** (沈): Often associated with the Shen River.} \\
\texttt{*   **Zhu** (朱): Historically linked to the Zhu River.} \\
\texttt{*   **Xia** (夏): Associated with the Xia River.} \\
\texttt{*   **Zhuang** (庄): Can be related to river settlements.} \\
\texttt{*   **Zhou** (周): While meaning "circle," it is phonetically similar to "Jiang" and historically linked to the Zhou River.} \\
\texttt{*   **Liu** (刘): Phonetically similar and historically linked to the Liu River.} \\
\texttt{*   **Zhu** (祝): Sometimes linked to river deities in ancient contexts.} \\
\texttt{*   **Zhu** (祝): In some contexts, associated with the Zhu River.} \\
\texttt{*   **Zhu** (祝): In some contexts, associated with the Zhu River.} \\
\\ 
\texttt{*Note: The most direct phonetic and semantic cousins to "Jiang" are surnames like **Song** (宋), **Shen** (沈), and **Zhu** (朱), which all carry historical ties to specific rivers.*} \\
            }
        }
    }
    & & \\ \\

    \theutterance \stepcounter{utterance}  
    & & & \multicolumn{2}{p{0.3\linewidth}}{
        \cellcolor[rgb]{0.9,0.9,0.9}{
            \makecell[{{p{\linewidth}}}]{
                \texttt{\tiny{[GM$|$GM]}}
                {'combination:two_responses': {}}
            }
        }
    }
    & & \\ \\

    \theutterance \stepcounter{utterance}  
    & & & \multicolumn{2}{p{0.3\linewidth}}{
        \cellcolor[rgb]{0.9,0.9,0.9}{
            \makecell[{{p{\linewidth}}}]{
                \texttt{\tiny{[GM$|$GM]}}
                \texttt{game\_result = WIN} \\
            }
        }
    }
    & & \\ \\

\end{supertabular}
}

\end{document}
